\section*{अध्याय \ref{ch:b1:comparative-planetology} --- तुलनात्मक ग्रह विज्ञान}

\begin{solution}{exo:b1:comparative-planetology:1}
$V = \tfrac43\pi(\qty{6.9911e7}{m})^3 = \qty{1.4313e24}{m^3}$ और
$\bar\rho = \num{1898.13e24}/\num{1.4313e24} = \qty{1326}{kg.m^{-3}}$, जल के घनत्व का केवल
1.33 गुना, यद्यपि बृहस्पति पृथ्वी से 318 गुना भारी है: वह मुख्यतः सबसे हल्के तत्वों से बना है।
\end{solution}

\begin{solution}{exo:b1:comparative-planetology:2}
बुध (त्रिज्या \qty{2440}{km}, \qty{5429}{kg.m^{-3}}) और शुक्र
(\qty{6052}{km}, \qty{5243}{kg.m^{-3}}) छोटे हैं और चट्टान से सघन हैं: \omterm{def:b1:comparative-planetology:planet-types}{शैल-ग्रह}, चट्टान और लोहे
के। शनि (\qty{60268}{km}, \qty{687}{kg.m^{-3}}) विशाल है और पानी से हल्का है: एक \omterm{def:b1:comparative-planetology:planet-types}{गैस दानव}, मुख्यतः
हाइड्रोजन और हीलियम। अरुण (\qty{25559}{km}, \qty{1270}{kg.m^{-3}}) एक दानव है पर शनि से कहीं छोटा
और लगभग दुगना सघन: एक \omterm{def:b1:comparative-planetology:planet-types}{हिम दानव}, जिसके जल, मीथेन और अमोनिया हाइड्रोजन से भारी हैं।
\end{solution}

\begin{solution}{exo:b1:comparative-planetology:3}
$T_\mathrm{eq} = (2601.3\times0.23/(4\times\num{5.6704e-8}))^{1/4} =
(\num{2.638e9})^{1/4} = \qty{226.6}{K}$। शुक्र को पृथ्वी से 1.9 गुना अधिक सूर्यप्रकाश मिलता है, पर उसके
बादल उसका \qty{77}{\percent} परावर्तित करते हैं: वह केवल $2601.3\times0.23 = \qty{598}{W.m^{-2}}$ सोखता
है, जबकि पृथ्वी के लिए $1361\times0.706 = \qty{961}{W.m^{-2}}$।
\end{solution}

\begin{solution}{exo:b1:comparative-planetology:4}
सबसे एकसमान से सबसे केंद्रित तक: चंद्रमा, मंगल, पृथ्वी, बृहस्पति, शनि। केंद्र पर एक बिंदु-द्रव्यमान का
$I = 0$ है, अतः $C = 0$।
\end{solution}

\begin{solution}{exo:b1:comparative-planetology:5}
$T = 15.945421\times86400 = \qty{1377684}{s}$;
$GM = 4\pi^2(\num{1.22187e9})^3/\num{1377684}^2 = \qty{3.794e16}{m^3.s^{-2}}$;
$M = \num{3.794e16}/\num{6.6743e-11} = \qty{5.685e26}{kg}$, \qty{568.32e24}{kg} के \qty{0.03}{\percent} के भीतर।
\end{solution}

\begin{solution}{exo:b1:comparative-planetology:6}
$S = 1361\times(149.6/228.0)^2 = \qty{585.9}{W.m^{-2}}$ मंगल पर और
$1361\times(149.6/778.5)^2 = \qty{50.26}{W.m^{-2}}$ बृहस्पति पर, सारणी के अनुरूप (मंगल पर छोटा अंतर माध्य
दूरियों के पूर्णांकन से आता है)।
\end{solution}

\begin{solution}{exo:b1:comparative-planetology:7}
$S \propto d^{-2}$ और $T_\mathrm{eq} \propto S^{1/4}$, अतः $T_\mathrm{eq}
\propto d^{-1/2}$। बृहस्पति के लिए, $T_\mathrm{eq} = (50.26\times0.657/(4\times
\num{5.6704e-8}))^{1/4} = \qty{109.8}{K}$। जो ग्रह सोखने से अधिक विकिरित करता है, उसका भीतरी स्रोत
अवश्य होगा: बृहस्पति अपने निर्माण से संचित, या धीरे संकुचित होते समय उत्पन्न ऊष्मा मुक्त कर रहा होगा।
\end{solution}

\begin{solution}{exo:b1:comparative-planetology:8}
शुक्र: $737 - 226.6 = \qty{510}{K}$, $T_\mathrm{eq}$ का \qty{225}{\percent}।
पृथ्वी: $T_\mathrm{eq} = (1361\times0.706/(4\sigma))^{1/4} = \qty{255.1}{K}$,
$288 - 255.1 = \qty{33}{K}$, \qty{13}{\percent}। मंगल: $214 - 209.8 =
\qty{4}{K}$, \qty{2}{\percent}।
\end{solution}

\begin{solution}{exo:b1:comparative-planetology:9}
$\mathrm dm = 4\pi\rho r^2\,\mathrm dr$ के साथ: $I = \int_0^R\tfrac23r^2\,4\pi
\rho r^2\,\mathrm dr = \tfrac{8\pi}{15}\rho R^5$, और $MR^2 = \tfrac43\pi\rho
R^5$, अतः $C = \tfrac{8}{15}\cdot\tfrac34 = \tfrac25$। एक खोखले कोश का सारा द्रव्यमान $r = R$ पर
होता है: $I = \tfrac23MR^2$, $C = \tfrac23$।
\end{solution}

\begin{solution}{exo:b1:comparative-planetology:10}
$\tfrac52C\bar\rho = 2.5\times0.35\times5429 = 4750.4$;
$x^2 = (4750.4 - 3300)/(5429 - 3300) = 1450.4/2129 = 0.6812$, $x = 0.8254$,
क्रोड-त्रिज्या $0.8254\times2439.7 = \qty{2014}{km}$। $x^3 = 0.5623$ और
$\rho_c = 3300 + 2129/0.5623 = \qty{7086}{kg.m^{-3}}$। क्रोड में
$0.5623\times7086/5429 = \qty{73}{\percent}$ द्रव्यमान है। बुध के क्रोड की त्रिज्या अपनी त्रिज्या का
\qty{83}{\percent} तक पहुँचती है, जबकि पृथ्वी के लिए $3485/6371 =
\qty{55}{\percent}$ है: बुध अधिकतर क्रोड है, जिसके ऊपर एक पतला शैलीय प्रावार है।
\end{solution}

\begin{solution}{exo:b1:comparative-planetology:11}
(क) फलन $t\mapsto t^4$ उत्तल है, अतः
$\bigl(\tfrac{T_1 + T_2}{2}\bigr)^4 \le \tfrac{T_1^4 + T_2^4}{2} =
T_\mathrm{eq}^4$, और समता केवल तब जब $T_1 = T_2$। (ख) $T_1^4 =
2\times270^4 - 100^4 = \num{1.0529e10}$, $T_1 = \qty{320}{K}$, और माध्य तापमान
$(320 + 100)/2 = \qty{210}{K}$ है, जो $T_\mathrm{eq}$ से \qty{60}{K} नीचे।
\end{solution}

\begin{solution}{exo:b1:comparative-planetology:12}
$m = 4\pi R^2p/g = 4\pi\times(\num{6.0518e6})^2\times\num{9.2e6}/8.9 =
\num{4.602e14}\times\num{1.034e6} = \qty{4.76e20}{kg}$: पृथ्वी के वायुमंडल के द्रव्यमान का 93 गुना, और
शुक्र के द्रव्यमान का \num{9.8e-5}। पृथ्वी में तुलनीय मात्रा में कार्बन है, पर वायु में नहीं, बल्कि
कार्बोनेट चट्टानों (चूना-पत्थर) में और महासागर में, जहाँ द्रव जल ने उसे संचित रहने दिया।
\end{solution}

\begin{solution}{pb:b1:comparative-planetology:1}
\textbf{1.} $T = 0.31891\times86400 = \qty{27553.8}{s}$।

\textbf{2.} $GM = 4\pi^2\times(\num{9.378e6})^3/\num{27553.8}^2 =
\qty{4.289e13}{m^3.s^{-2}}$।

\textbf{3.} $M = \num{4.289e13}/\num{6.6743e-11} = \qty{6.426e23}{kg}$,
सारणीबद्ध \qty{6.4169e23}{kg} के \qty{0.14}{\percent} के भीतर।

\textbf{4.} $\num{1.06e16}/\num{6.426e23} = \num{1.6e-8}$: फ़ोबोस नगण्य है।

\textbf{5.} $v = 2\pi a/T = 2\pi\times\num{9.378e6}/27553.8 =
\qty{2138}{m.s^{-1}}$, लगभग \qty{2.1}{km.s^{-1}}।

\textbf{6.} $T = \qty{7.65}{h}$, मंगल के घूर्णन से 3.2 गुना छोटा: फ़ोबोस ज़मीन के अपने नीचे घूमने से
तेज़ परिक्रमा करता है, इसलिए वह घूर्णन को पीछे छोड़ता है। मंगल से देखने पर वह पश्चिम में उगता और
पूर्व में अस्त होता है।

\textbf{7.} $V = \tfrac43\pi(\num{3.3895e6})^3 = \qty{1.6312e20}{m^3}$।

\textbf{8.} $\bar\rho = \num{6.426e23}/\num{1.6312e20} = 3939$, अर्थात्
\qty{3940}{kg.m^{-3}}।

\textbf{9.} $3940/5514 = 0.715$। पृथ्वी, जो बड़ी और भारी है, अपने भीतर को अधिक प्रबलता से संपीडित
करती है; और उसका लौह-क्रोड ग्रह का बड़ा भाग है।

\textbf{10.} $g = GM/R^2 = \num{4.289e13}/(\num{3.3895e6})^2 =
\qty{3.73}{m.s^{-2}}$, पृथ्वी के \qty{38}{\percent}।

\textbf{11.} $v_\mathrm{esc} = \sqrt{2GM/R} = \sqrt{2\times\num{4.289e13}/
\num{3.3895e6}} = \qty{5.03}{km.s^{-1}}$।

\textbf{12.} $T_\mathrm{eq} = (586.2\times0.75/(4\times\num{5.6704e-8}))^{1/4}
= \qty{209.8}{K}$।

\textbf{13.} $214 - 209.8 = \qty{4}{K}$: एक दुर्बल हरितगृह तापन।

\textbf{14.} $m = p/g = 636/3.733 = \qty{170}{kg.m^{-2}}$; पृथ्वी के लिए
$101400/9.82 = \qty{10326}{kg.m^{-2}}$, 61 गुना अधिक।

\textbf{15.} $4\pi R^2\times170.4 = \num{1.444e14}\times170.4 =
\qty{2.46e16}{kg}$।

\textbf{16.} हर वर्ग मीटर के ऊपर पृथ्वी की तुलना में साठ गुना कम गैस है, इसलिए अवरक्त अवशोषण, यद्यपि
कार्बन डाइऑक्साइड द्वारा, दुर्बल है; और मंगल की वायु में बहुत थोड़ा जल-वाष्प है, जो दूसरा प्रमुख
अवशोषक है।

\textbf{17.} एकसमान मंगल का $C = 0.4$ होता; मापा गया 0.366 दिखाता है कि उसका घनत्व केंद्र की ओर बढ़ता
है: मंगल का एक सघन क्रोड है।

\textbf{18.} $\bar\rho = \rho_m + (\rho_c - \rho_m)x^3$।

\textbf{19.} $C = \tfrac25[\rho_m + (\rho_c - \rho_m)x^5]/\bar\rho$।

\textbf{20.} 18 और 19 से, $(\rho_c - \rho_m)x^3 = \bar\rho - \rho_m$ और
$(\rho_c - \rho_m)x^5 = \tfrac52C\bar\rho - \rho_m$; भाग देने पर $x^2$ मिलता है।

\textbf{21.} $\tfrac52C\bar\rho = 2.5\times0.366\times3940 = 3605.1$;
$x^2 = (3605.1 - 3500)/(3940 - 3500) = 105.1/440 = 0.2389$, $x = 0.4887$।
$x^3 = 0.1167$ और $\rho_c = 3500 + 440/0.1167 = \qty{7269}{kg.m^{-3}}$।

\textbf{22.} $x^3\rho_c/\bar\rho = 0.1167\times7269/3940 = 0.215$, अर्थात् \qty{21.5}{\percent} द्रव्यमान।

\textbf{23.} $\rho_m = 3450$ के साथ: $x^2 = 155.1/490 = 0.3165$, $x = 0.5626$,
क्रोड-त्रिज्या \qty{1907}{km}, $\rho_c = \qty{6202}{kg.m^{-3}}$। $\rho_m =
3550$ के साथ: $x^2 = 55.1/390 = 0.1413$, $x = 0.3759$, \qty{1274}{km},
$\rho_c = \qty{10894}{kg.m^{-3}}$। अनुमानित प्रावार घनत्व में \qty{1.4}{\percent} का परिवर्तन क्रोड-त्रिज्या को
कई सौ किलोमीटर हिला देता है: माध्य घनत्व और $C$ अकेले मंगल के क्रोड को सूक्ष्मता से तय नहीं कर सकते, इसीलिए
वहाँ भूकंपमापी भेजे गए।

\textbf{24.} $0.4887\times3389.5 = \qty{1657}{km}$: मॉडल से मंगल के क्रोड की त्रिज्या लगभग
\textbf{\qty[text-series-to-math=true]{1660}{km}} है, जो हाल की भूकंपीय सीमा \qty{1650}{km} से \qty{1675}{km}
के भीतर है और ग्रहीय सारणी के \qty{1830}{km} से काफ़ी नीचे है --- एक ऐसी सहमति जो प्रश्न 23 दिखाता है कि
$\rho_m$ के चुनाव में आंशिक रूप से संयोग है।
\end{solution}
